
\section{REINFORCE Gradient Check}
\label{sec:appendix_gradient}

The directional correctness of the \reinforce{} gradient estimator was validated against a symmetric finite-difference approximation~\cite{williams1992simple}.

\subsection{Method}
For a fixed policy parameterization $\theta$ and random seed, we computed:
(i)~the \reinforce{} gradient estimate $\hat{g}_{\mathrm{RF}}$, and
(ii)~the finite-difference gradient $\hat{g}_{\mathrm{FD}}$ using symmetric perturbations.
Three metrics were evaluated: cosine similarity, relative $\ell_2$ error, and bias estimate.

\subsection{Results}

\begin{table}[htbp]
\centering
\small
\caption{Gradient check results.}
\label{tab:appendix_gradient}
\renewcommand{\arraystretch}{1.25}
\begin{tabular}{@{\hspace{6pt}}l @{\hspace{14pt}} l @{\hspace{14pt}} p{5.5cm}@{\hspace{6pt}}}
\toprule
\textbf{Metric} & \textbf{Value} & \textbf{Interpretation} \\
\midrule
Cosine similarity & 1.0000 & Perfect directional agreement \\
Relative error    & 32.61\% & Magnitude mismatch \\
Bias estimate     & $1.19 \times 10^{-3}$ & Negligible systematic bias \\
Variance estimate & $4.18 \times 10^{-9}$ & Near-zero \\
\bottomrule
\end{tabular}
\end{table}

The cosine similarity of $1.0000$ confirms exact directional agreement.
The 32.6\% relative magnitude error is attributable to sample-size differences between the two estimators.
For policy gradient training, directional correctness is the critical requirement; magnitude affects learning-rate tuning but not the optimization direction.

\begin{figure}[htbp!]
\centering
\includegraphics[width=0.70\textwidth, height=0.35\textheight, keepaspectratio]{gradient_scatter.png}
\caption{Parameter-wise scatter: \reinforce{} gradient components vs.\ finite-difference gradients.  Near-perfect alignment along the diagonal confirms directional correctness.}
\label{fig:appendix_gradient}
\end{figure}

% Paths removed for double-blind review.


